\section{Introduction}
\label{sec:intro}

A deployed classifier can fail after deployment for different reasons, and the reasons call for different responses. If only the input distribution moves while the conditional law of the label stays the same (covariate shift), the model's posterior remains correct and the usual response is to keep the model and keep watching. If the class prevalence moves while the class-conditional input laws stay fixed (label shift), the posterior is repaired by a prior correction that needs no new training \citep{saerens2002,lipton2018,alexandari2020}. If the relation between inputs and labels changes in a way that no prior correction repairs (concept shift), the model must be retrained or must abstain. A monitor that only reports that ``something changed'' leaves this decision to guesswork. A monitor that names a type without any error control can trigger an unnecessary retraining, or an inappropriate prior correction, with no bound on how often this happens.

Two further constraints make the problem hard in practice. Labels are sparse: only a small random fraction of outcomes is ever observed, so evidence about $P(Y\mid X)$ accumulates slowly, while unlabelled inputs arrive at every step. And monitoring is continuous: the operator looks at the monitor at data-dependent times and acts on its first alarm, so false-alarm control must hold under optional stopping and, for change detection, should not degrade when the change arrives late.

This paper develops \emph{Typed Shift Sequential Diagnosis} (TSSD) for this setting. TSSD runs three witnesses on one false-alarm budget. Two are paired, source-referenced e-processes: an orthogonalised kernel witness for inputs and a gated prevalence witness that also uses unlabelled cases. The third is new: a \emph{posterior-validity} witness, a universal-inference e-process \citep{wasserman2020} whose composite null is the set of all target posteriors that a prior and tempering correction of the frozen model can reach, \emph{with a label shift allowed at an unknown time and with no restriction on the inputs}. Its crossing is therefore a certified statement that prior correction is insufficient, the decision that separates retraining from recalibration. Every witness is evaluated in two modes from the same increments: a test-martingale mode with ``ever-alarm'' control, and a Shiryaev--Roberts e-detector mode \citep{shin2024edetectors} that restarts implicitly and controls the average run length.

\paragraph{Contributions.}
\begin{enumerate}[leftmargin=1.4em,itemsep=0.15em]
\item \emph{Action-typed formulation with a certified concept claim.} We define nested monitoring hypotheses aligned with the keep/recalibrate/retrain decision. We also construct a posterior-validity e-process whose false ``concept'' claims are controlled at level $\alpha/3$ under arbitrary covariate shift and under a label shift at an unknown time (Theorem~\ref{thm:concept}). The required change-robust likelihood maximisation is computed exactly from cached suffix suprema. A certified concavity bound makes the denominator valid even if the optimiser stops early.
\item \emph{Change-detection guarantees.} A Shiryaev--Roberts version of all three witnesses gives a typed alarm with $\E[\tau]\ge A/3$ and $\Pr(\tau\le m)\le3m/A$ under the global null (Theorem~\ref{thm:sr}). Empirically, its detection delay does not grow with the change time, whereas the test-martingale delay does.
\item \emph{Exact validity with finite source pools.} Single-use references drawn without replacement make the guarantees exact for real-data pools (Proposition~\ref{prop:pool}).
\item \emph{A broad empirical comparison.} We compare TSSD on identical streams with ten re-implemented detectors that span the literature from DDM and ADWIN to WATCH, and with a Bonferroni ensemble of specialist detectors. The comparison covers six pure and mixed shift types, graded magnitudes, four label budgets, four change times and eleven ablations, and we report both native thresholds and thresholds matched to a common empirical false-alarm rate. The real-data evaluation uses three datasets: Digits under an exact-validity protocol, natural geographic shift in Covertype, and chronological drift in Electricity. For the two natural shifts the shift type is checked against an offline full-label oracle.
\end{enumerate}

\paragraph{What TSSD is not.} TSSD is not the fastest detector of every individual shift type. A specialist detector tuned to one alternative (for instance Page--Hinkley on the log loss for concept shift) detects that alternative sooner, and we report this. TSSD's contribution is a calibrated typed state with a guarantee on its most consequential claim. In our experiments, the obvious alternative, an ensemble of specialists typed by which detector fired, assigns the wrong type to most covariate shifts and makes spurious concept claims under mixed covariate--label shift (Section~\ref{sec:typing}). TSSD does not adapt the model, and its distinction between covariate and label shift is heuristic and is evaluated empirically.
